% =====================================================================
%  ЗУРГАА. САНХҮҮГИЙН ТӨЛӨВЛӨГӨӨ
% =====================================================================
\chapter{САНХҮҮГИЙН ТӨЛӨВЛӨГӨӨ}

\section{Үйлдвэрлэлийн зардлын тооцоо}

Төслийн хэвийн ажиллагаанд шилжсэний дараах сар бүрийн тогтмол
зардал (Хүснэгт~\ref{tab:ch6-zardal}).

\begin{table}[h]
  \centering
  \caption{Сар бүрийн үйл ажиллагааны зардал /төгрөг/}
  \label{tab:ch6-zardal}
  \small
  \begin{tabular}{@{}p{0.8cm}p{7.6cm}r@{}}
    \toprule
    \textbf{№} & \textbf{Зардлын төрөл} & \textbf{Сарын дүн} \\
    \midrule
    1 & Cloud сервер, CDN, хадгаламж & 300{,}000 \\
    2 & Техник хангамжийн дэмжлэг, засвар үйлчилгээ & 250{,}000 \\
    3 & Ажилтны цалин (үйлчилгээний баг, 3 хүн) & 2{,}500{,}000 \\
    4 & Даатгал, НДШ & 315{,}000 \\
    5 & Маркетинг, хэрэглэгч татах зардал & 400{,}000 \\
    6 & Кофе шопын стикер солих, нэмэлт тоног төхөөрөмж & 100{,}000 \\
    7 & Бусад (офис, тээвэр, нөөц) & 235{,}000 \\
    \midrule
    & \textbf{Нийт} & \textbf{4{,}100{,}000} \\
    \bottomrule
  \end{tabular}
\end{table}

\section{Үйлдвэрлэлийн орлогын тооцоо (Борлуулалтын төлөвлөгөө)}

Орлогын гол эх үүсвэр нь кофе шопуудад зарах нэг удаагийн багц үнэ юм.
Урьдчилсан тооцоог Хүснэгт~\ref{tab:ch6-orlogo}~—т харуулав.

\begin{table}[h]
  \centering
  \caption{Борлуулалтын урьдчилсан тооцоо /төгрөг/}
  \label{tab:ch6-orlogo}
  \small
  \begin{tabular}{@{}p{0.8cm}p{4.4cm}rrr@{}}
    \toprule
    \textbf{№} & \textbf{Үзүүлэлт} & \textbf{1-р жил} & \textbf{2-р жил} &
    \textbf{3-р жил} \\ \midrule
    1 & Борлуулсан кофе шоп (Багц) & 15 & 25 & 40 \\
    2 & Багцын дундаж үнэ & 3{,}300{,}000 & 3{,}500{,}000 & 3{,}700{,}000 \\
    3 & Багцын борлуулалтын орлого & 49{,}500{,}000 & 87{,}500{,}000 & 148{,}000{,}000 \\
    4 & Нэмэлт стикер, жилийн дэмжлэгийн орлого & 3{,}000{,}000 &
        10{,}000{,}000 & 15{,}000{,}000 \\
    5 & \textbf{Нийт орлого} & \textbf{52{,}500{,}000} &
        \textbf{97{,}500{,}000} & \textbf{163{,}000{,}000} \\
    \bottomrule
  \end{tabular}
\end{table}

\section{Орлого, зардлын нэгдсэн тооцоолол}

\begin{table}[h]
  \centering
  \caption{Орлого, зардлын нэгдсэн тооцоолол /төгрөг/}
  \label{tab:ch6-negtetsen}
  \small
  \begin{tabular}{@{}p{0.6cm}p{6.4cm}rrr@{}}
    \toprule
    \textbf{№} & \textbf{Үзүүлэлт} & \textbf{1-р жил} & \textbf{2-р жил} &
    \textbf{3-р жил} \\ \midrule
    1 & Нийт орлого & 52{,}500{,}000 & 97{,}500{,}000 & 163{,}000{,}000 \\
    2 & Үйл ажиллагааны зардал (4.1\,сая x 12) & 49{,}200{,}000 &
        60{,}000{,}000 & 80{,}000{,}000 \\
    3 & Өгөгдөл хадгалах, API-ийн нэмэлт зардал & 3{,}000{,}000 &
        6{,}000{,}000 & 9{,}000{,}000 \\
    4 & ТТӨА (өмнөх) & 300{,}000 & 31{,}500{,}000 & 74{,}000{,}000 \\
    5 & Корпорацийн албан татвар (10\%) & 30{,}000 & 3{,}150{,}000 & 7{,}400{,}000 \\
    6 & \textbf{Цэвэр ашиг} & \textbf{270{,}000} &
        \textbf{28{,}350{,}000} & \textbf{66{,}600{,}000} \\
    \bottomrule
  \end{tabular}
\end{table}

\noindent 1-р жилд хөгжүүлэлтийн зардлыг хөрөнгө оруулалтад тусгаж
ажиллагааны зардлыг орлогоор бүрэн хамруулах түвшинд ажиллахаар
төлөвлөв. 2-р жилээс тогтвортой ашигтай ажиллана.

\section{Хөрөнгө оруулалтын тооцоо}

Төслийг хэрэгжүүлэхэд шаардагдах нийт хөрөнгө оруулалтыг
Хүснэгт~\ref{tab:ch6-hurungoo}~—т тусгав.

\begin{table}[h]
  \centering
  \caption{Хөрөнгө оруулалтын нийт зардал /төгрөг/}
  \label{tab:ch6-hurungoo}
  \small
  \begin{tabular}{@{}p{0.8cm}p{8.2cm}r@{}}
    \toprule
    \textbf{№} & \textbf{Хөрөнгө оруулалтын төрөл} & \textbf{Дүн} \\
    \midrule
    1 & NFC стикер, үйлчилгээний төхөөрөмж (300 ширээ) & 2{,}000{,}000 \\
    2 & Хөгжүүлэлтийн тоног төхөөрөмж (компьютер, тест утас) & 5{,}000{,}000 \\
    3 & Програм хангамжийн хөгжүүлэлт (6 сарын багийн зардал) & 34{,}185{,}360 \\
    4 & Cloud сервер, домайн, SSL (1 жил) & 2{,}400{,}000 \\
    5 & Төлбөрийн интеграци, лиценз, нотолгоо & 1{,}500{,}000 \\
    6 & Маркетинг, нэвтрүүлгийн зардал & 3{,}000{,}000 \\
    7 & Ажлын байрны зардал (6 сар) & 2{,}400{,}000 \\
    8 & Хөгжлийн нөөц (10\%) & 5{,}048{,}536 \\
    \midrule
    & \textbf{Нийт дүн} & \textbf{55{,}533{,}896} \\
    \bottomrule
  \end{tabular}
\end{table}

\noindent Тэмдэглэл: Хөгжүүлэлтийн зардлыг санхүүгийн тооцоонд
бүрэн хэмжээгээр оруулсан бөгөөд дараагийн үе шатанд бүтээгдэхүүний
хөгжүүлэлт орлогоороо санхүүжигдэнэ.

\section{Хөрөнгө оруулалтын үр ашиг}

\subsection{Хөрөнгө оруулалтын өгөөж (ROI)}

\begin{equation}
  ROI = \frac{\text{3 жилийн нийт цэвэр ашиг}}{\text{Нийт хөрөнгө оруулалт}}
      = \frac{95{,}220{,}000}{55{,}533{,}896} \approx 1.71 \; (171\%)
\end{equation}

\noindent Энэ нь нэгж хөрөнгө оруулалт тутамд 3 жилийн хугацаанд
1.71 төгрөгийн цэвэр ашиг олохоор байна гэсэн үг.

\subsection{Хөрөнгө оруулалтыг нөхөх хугацаа}

2-р жилээс эхлэн сар бүр 2{,}362{,}500\textcurrency{} (28{,}350{,}000 / 12)
цэвэр ашиг олохоор тооцсон үзүүлэлтээр, нийт хөрөнгө оруулалтын
1-р жилийн төгсгөлд үлдсэн хэсгийг нөхөхөд:

\begin{equation}
  \text{Нөхөх хугацаа} \approx 12 + \frac{55{,}263{,}896}{2{,}362{,}500}
  \approx 35 \text{ сар}
\end{equation}

\noindent Өөрөөр хэлбэл, төсөл эхэлснээс хойш ойролцоогоор 35 сарын
дотор хөрөнгө оруулалтаа бүрэн нөхөж, ажиллагааны зардлаа орлогоороо
бүрэн хамруулах түвшинд хүрнэ (борлуулалтын эрчмээс хамаарна).

\noindent Хөгжүүлэлтийн зардлыг (34{,}185{,}360\textcurrency{})
хуримтлагдсан цэвэр ашигаар нөхөхийг тооцвол: 1-р жилийн 270{,}000\textcurrency{}
ба 2-р жилийн сар бүрийн 2{,}362{,}500\textcurrency{} ашигаар
хөгжүүлэлтийн зардлыг ойролцоогоор 26 сарын дотор (3-р жилийн
дунд үе) бүрэн нөхөнө. Ингэснээр 3-р жилийн цэвэр ашиг
(66{,}600{,}000\textcurrency{}) бүхэлдээ хөрөнгө оруулалтын өгөөж болно.

\section{Зээлийн эргэн төлөлтийн тооцоо}

Хөрөнгө оруулалтын 50\% (27{,}766{,}948\textcurrency{}) нь 3 жилийн
хугацаатай, жилд 12\%-ийн хүүтэй зээлээр хангагдана
(Хүснэгт~\ref{tab:ch6-zeel}).

\begin{table}[h]
  \centering
  \caption{Төслийн зээлийн төлөлтийн график /төгрөг/}
  \label{tab:ch6-zeel}
  \small
  \begin{tabular}{@{}p{1.4cm}p{2.2cm}rrrr@{}}
    \toprule
    \textbf{Он} & \textbf{Төлөх сар} & \textbf{Үндсэн төлбөр} &
    \textbf{Зээлийн хүү} & \textbf{Нийт төлбөр} & \textbf{Зээлийн
    үлдэгдэл} \\ \midrule
    1 & 12-р сар & 9{,}255{,}649 & 3{,}332{,}034 & 12{,}587{,}683 &
    18{,}511{,}299 \\
    2 & 24-р сар & 9{,}255{,}649 & 2{,}221{,}356 & 11{,}477{,}005 &
    9{,}255{,}650 \\
    3 & 36-р сар & 9{,}255{,}650 & 1{,}110{,}678 & 10{,}366{,}328 & 0 \\
    \midrule
    & \textbf{Нийт} & \textbf{27{,}766{,}948} & \textbf{6{,}664{,}068} &
    \textbf{34{,}431{,}016} & \\
    \bottomrule
  \end{tabular}
\end{table}
